\cgeentry{cge_dynamic}{Solve a dynamic run, recursive or with perfect foresight}

\cgesec{Description}

\cmd{cge_dynamic} solves the calibrated model for a sequence of periods. The
first period is the calibrated base; each later period is an equilibrium of
the static model given stocks and exogenous values that are updated from the
solution of the period before. Agents do not look ahead, unless the run has
perfect foresight (\opt{fores}{ight}, from version 0.10.2): the periods are
then solved together, and user equations may refer to later periods with
\cmd{lead()}. Each period is
stored as a simulation named after the run and the period, such as
\cmd{base2025}, so that \manref{cge_report}, \manref{cge_results},
\manref{cge_dynresults}, \manref{cge_save} and \manref{cge_gams} work with it.
The run uses the current closure and starts from the base; the parameters and
fixed values that the model had before the run are restored when it ends.
If a period cannot be solved, or the \opt{shock}{s()} program fails, the run
stops and nothing is stored: the model, its shocks and its stored simulations
are as they were before the run.

\cgesec{Quick start}

A baseline from 2018 to 2030 in which capital accumulates, labour grows at 3
per cent, the population at 2.5 per cent, and efficiency at the rate that
gives 5 per cent growth of real value added
\begin{stcode}
cge_dynamic, name(base) periods(2018/2030) capital(cap) depreciation(0.05) rental(0.15) ///
    growth(lab:0.03) population(0.025) gdpgrowth(0.05)
\end{stcode}
A scenario with the same settings, the baseline's efficiency path and the
shocks applied by program \cmd{policy}
\begin{stcode}
cge_dynamic, name(policy) periods(2018/2030) capital(cap) depreciation(0.05) rental(0.15) ///
    growth(lab:0.03) population(0.025) baseline(base) shocks(policy)
\end{stcode}
Ten periods with efficiency growth of 1 per cent in every activity and
government flows that grow at 3 per cent
\begin{stcode}
cge_dynamic, name(run) periods(0/10) tfp(0.01) govgrowth(0.03)
\end{stcode}
Labour that loses the deaths at working ages, in a model calibrated with a
death rate
\begin{stcode}
cge_dynamic, name(hb) periods(0/10) growth(lab:0.02) population(0.02) ///
    deathfactors(lab) deathshare(0.3) workingage(0.55)
\end{stcode}
The periods solved together, with perfect foresight, and leads beyond the last
period that continue its growth
\begin{stcode}
cge_dynamic, name(fs) periods(2018/2030) capital(cap) depreciation(0.05) rental(0.10) ///
    growth(lab:0.02) foresight terminal(growth)
\end{stcode}

\cgesec{Syntax}

\begin{stsyntax}
\cmdb{cge_dynamic}\texttt{,} \optarg{}{name}{name} \optarg{per}{iods}{numlist}
[\argu{options}]
\end{stsyntax}

\begin{opttable}
\optgroup{Main}
\optarg{}{name}{name} & name of the run; required\\
\optarg{per}{iods}{numlist} & periods, the base first; required\\
\optarg{shock}{s}{program} & program that applies the shocks of each period\\
\optgroup{Capital}
\optarg{cap}{ital}{factors} & factors that accumulate\\
\optarg{dep}{reciation}{numlist} & depreciation rates of the capital factors\\
\optarg{rent}{al}{numlist} & base gross rental rates of the capital factors\\
\optarg{sens}{itivity}{\#} & response of the allocation of new capital to rental rates; default \texttt{1}\\
\optgroup{Growth}
\optarg{grow}{th}{factor\texttt{:}rate \ldots} & growth rates of other factors\\
\optarg{pop}{ulation}{\#} & population growth rate; default \texttt{0}\\
\optarg{gov}{growth}{\#} & growth rate of government consumption and transfers; default \texttt{0}\\
\optarg{for}{eigngrowth}{\#} & growth rate of flows in foreign currency; default \texttt{0}\\
\optgroup{Efficiency}
\optarg{}{tfp}{numlist} & growth rate of efficiency, one or one per activity\\
\optarg{gdp}{growth}{numlist} & growth targets for real value added, one or one per period after the base\\
\optarg{base}{line}{name} & reuse the efficiency path of run \argu{name}\\
\optgroup{Deaths}
\optarg{deathf}{actors}{factors} & labour factors that lose workers to deaths\\
\optarg{deaths}{hare}{\#} & share of deaths at working ages; default \texttt{0}\\
\optarg{work}{ingage}{\#} & working-age share of the population; default \texttt{1}\\
\optgroup{Perfect foresight}
\opt{fores}{ight} & solve the periods together\\
\optarg{term}{inal}{constant|growth|base} & value of a lead beyond the last period; default \cmd{constant}\\
\optarg{}{start}{recursive|base} & start of Newton's method; default \cmd{recursive}\\
\optgroup{Solver}
\optarg{tol}{erance}{\#} & tolerance of each period's solve; default \texttt{1e-12}\\
\optarg{iter}{ate}{\#} & maximum number of Newton iterations of each solve; default \texttt{100}\\
\end{opttable}

Growth rates are per period and are zero unless they are given. At most one
of \cmd{tfp()}, \cmd{gdpgrowth()} and \cmd{baseline()} may be specified.

\cgesec{Options}

\begin{optdesc}
\item[\optarg{}{name}{name}] names the run. Each period is stored as the
simulation \argu{name} followed by the period, for example \cmd{base2025}, so
\argu{name} with each period must form a valid Stata name of at most 32
characters. Simulations of the same name are replaced, and a run solved again
replaces every period of the earlier run of that name.

\item[\optarg{per}{iods}{numlist}] lists the periods: at least two integers,
zero or positive, in ascending order, for example \cmd{periods(2018/2030)} or
\cmd{periods(0/5)}. The first is the base. Each step from one listed period to
the next is one period of the rates, whatever the difference between the two
numbers.

\item[\optarg{shock}{s}{program}] names a program that is called as
\argu{program period} before each period after the first, once stocks and
exogenous values have been updated and before efficiency is set. It may apply
any shock with \manref{cge_shock}, for example to cancel debt with
\cmd{cge_shock fdebt}. Shocks persist in later periods, so a shock with
\cmd{scale()} should be applied in one period only. If the program fails, the
run stops, and nothing is stored.

\item[\optarg{cap}{ital}{factors}] lists the factors that accumulate capital.
Each must be mobile or specific to each activity in the closure, not
unemployed or under a wage curve or minimum wage. \cmd{capital()} needs
\cmd{depreciation()} and \cmd{rental()}, which need it.

\item[\optarg{dep}{reciation}{numlist}] gives the depreciation rates of the
capital factors, at least 0 and less than 1: one rate for all, or one for each
factor in the order of \cmd{capital()}.

\item[\optarg{rent}{al}{numlist}] gives the base gross rental rates of the
capital factors, the rental per unit of capital at base prices, so that the
base capital stock is base capital income divided by the rate. One unit of
real investment adds \argu{rental} units of the factor. Rates are positive:
one for all, or one for each factor in the order of \cmd{capital()}.

\item[\optarg{sens}{itivity}{\#}] sets how strongly new capital specific to
each activity moves to activities whose rental rate is above the average; 0
allocates it in proportion to capital stocks. The default is 1; \# must be
zero or positive. It has no effect on mobile capital.

\item[\optarg{grow}{th}{factor\texttt{:}rate \ldots}] sets the growth rates of
factors that are not capital; factors that are not listed do not grow. With
\cmd{growth(lab_f:0.03 lab_i:0.025)}, for example, \cmd{lab_f} grows at 3 per
cent a period and \cmd{lab_i} at 2.5 per cent. The supply of an unemployed factor is
endogenous, so its rate must be zero, unless the factor has a labour force
(\manref{cge_calibrate}, \cmd{unemployment()}). The labour force of a factor
that has one grows at the factor's rate, whatever its closure, so that under
\cmd{unemployed}, \cmd{wagecurve} or \cmd{minwage} employment and
unemployment follow from labour demand.

\item[\optarg{pop}{ulation}{\#}] sets the population growth rate. Subsistence
consumption of LES demand, the base of the health link and the population
index grow at this rate. With a death rate (\mbox{\manref{cge_calibrate}},
\cmd{mortality()}), the rate is taken to hold at the base death rate:
population then grows at \cmd{population()} plus the base death rate less the
death rate of the period before.

\item[\optarg{gov}{growth}{\#}] sets the growth rate of government
consumption and of real transfers from the government to domestic
institutions, within the government and to the health insurance fund.

\item[\optarg{for}{eigngrowth}{\#}] sets the growth rate of flows in foreign
currency: transfers from abroad to domestic institutions, factors, the
government and the health insurance fund, transfers from the government
abroad, and foreign savings when the closure fixes them.

\item[\optarg{}{tfp}{numlist}] sets the growth rate of value-added efficiency:
one rate, used for every activity, or one for each activity in the order of
the model's activities.

\item[\optarg{gdp}{growth}{numlist}] calibrates a baseline to growth targets:
in each period, efficiency in every activity grows by the uniform factor for
which real value added at base prices grows at the target. Give one target,
used in every period, or one for each period after the first. The efficiency
factor of each period is displayed and stored with its simulation.

\item[\optarg{base}{line}{name}] reuses the efficiency factors of run
\argu{name}, which must have been solved with \cmd{gdpgrowth()} and must
include every period of this run, so that a policy scenario does not change
efficiency.

\item[\optarg{deathf}{actors}{factors}] lists the labour factors that lose
workers to deaths at working ages. They may not be capital factors, nor
unemployed factors without a labour force, and the model needs a death rate.
Each listed factor, and its labour force, then grows at its \cmd{growth()}
rate plus $\text{\cmd{deathshare()}}/\text{\cmd{workingage()}}$ times the base
death rate less the death rate of the period before.

\item[\optarg{deaths}{hare}{\#}] is the share of deaths at working ages,
between 0 and 1; the default is 0. It needs \cmd{deathfactors()}.

\item[\optarg{work}{ingage}{\#}] is the working-age share of the population,
above 0 and at most 1; the default is 1. It needs \cmd{deathfactors()}.

\item[\opt{fores}{ight}] (from version 0.10.2) solves the periods of the run
together, with perfect foresight (\secref{sec:foresight}). The updates that do
not depend on a solution and the shocks of each period are made first; capital
stocks, public debt, efficiency under growth targets and, with a death rate,
population and the factors that lose workers become unknowns of the stacked
system, with their updates as equations. User equations may then refer to
later periods, \cmd{lead(VAR)} and \cmd{lead(VAR(e), k)}, $k \le 4$; a run
whose user equations have leads needs \opt{fores}{ight}. Without leads the
run gives the recursive run. The \cmd{shocks()} program runs twice in each
period, the second time quietly, so that a shock to a capital or debt stock
is applied to the stock that the run finds; with \cmd{gdpgrowth()} efficiency
may not be shocked, nor, with a death rate, population or the factors in
\cmd{deathfactors()}, and the closure stays the same.

\item[\optarg{term}{inal}{constant|growth|base}] sets the value of a lead
beyond the last period $T$: that of the last period (\cmd{constant}, the
default), the growth of the last period continued, $\V{VAR}_T
(\V{VAR}_T/\V{VAR}_{T-1})^k$ (\cmd{growth}), or the base (\cmd{base}). An
equation of its own for the last period is written with
\cmd{ifthen(tlast,} \ldots\cmd{)} (\manref{cge_equation}).

\item[\optarg{}{start}{recursive|base}] sets the start of Newton's method on
the stacked system: a recursive run in which each period expects the later
ones to repeat the period before (\cmd{recursive}, the default), or the base
in every period (\cmd{base}). Both reach the same solution; the first is
usually closer.

\item[\optarg{tol}{erance}{\#}] and \optarg{iter}{ate}{\#} are passed to the
solver of each period, or of the stacked system with \opt{fores}{ight}; the
defaults are \texttt{1e-12} and 100.
\end{optdesc}

\cgesec{Remarks and examples}

A dynamic run strings static equilibria together. Before each period after
the first, \cmd{cge_dynamic} updates, in this order, capital, the supply of
other factors, the population, government and foreign flows and public debt;
it then calls the \cmd{shocks()} program and sets efficiency, and solves the
period from the solution of the period before. The updates follow the
recursive dynamic extension of the standard model (Thurlow 2004);
Methods and formulas below gives them, and \chapref{chap:dynamics} the
details. None of the rates has a default other than zero, and none is
estimated by the package: they are the user's assumptions about the economy.

\cgesubsec{Example 1: a baseline calibrated to a growth target}

The synthetic SAM is calibrated as in \manref{cge_solve}, and capital is made
specific to each activity, so that new capital is allocated among activities:
\stlog{dyn_setup}
A baseline of four periods after the base, in which capital depreciates at 5
per cent a period, labour and the population grow at 2.5 per cent, government
flows at 3 per cent and foreign flows at 2 per cent, and efficiency grows at
the rate that gives 4 per cent growth of real value added:
\stlog{dyn_base}
Period 0 is the base and needs no iteration. In each later period, the
efficiency factor is the uniform growth factor of value-added efficiency
that, given the growth of capital and labour, meets the target; it is found by
secant iterations, each of which solves the model. The run returns the names
of its simulations:
\stlog{dyn_return}
\manref{cge_report} with \cmd{run()} reports every period as a change from the
base:
\stlog{dyn_basereport}
Real value added grows at 4 per cent a period, 16.986 per cent in four
periods ($1.04^4 = 1.16986$). The SAM is synthetic: the run illustrates the
command and is not a projection for any economy.

\cgesubsec{Example 2: a policy scenario against the baseline}

A scenario removes every import tariff from period 2. The program
\cmd{tariff_from2} receives the period as its argument; the shock is set once
and stays in force in later periods. The scenario takes the baseline's
settings and, with \cmd{baseline(base)}, its efficiency factors:
\stlog{dyn_scenario}
The efficiency factors are those of the baseline, period by period, so that
efficiency does not respond to the policy. The shock is reported before
period 2, the first period that it changes. Periods of the two runs can be
compared side by side:
\stlog{dyn_compare}
Both columns are changes from the base, period 0. In period 4 real investment
is 3.5 per cent above the base in the scenario, against 18.9 per cent in the
baseline, and real value added grows less, 16.1 against 17.0 per cent; real
consumption and the equivalent variation of the household are higher. New
capital is proportional to real investment in the period before, so lower
investment means less capital in later periods. The differences
from the baseline, period by period and for every variable, are the dataset
of \manref{cge_dynresults}.

\cgesubsec{Example 3: population, deaths and labour}

With a death rate (\manref{cge_calibrate}, \cmd{mortality()} and
\cmd{mortelas()}), the deaths of one period change the population and labour
of the next. The health layout of the synthetic SAM is calibrated with a base
death rate of 10 per 1,000 and an elasticity of $-0.3$ with respect to health
services per head. In the run, 30 per cent of deaths are at working ages, and
55 per cent of the population is of working age:
\stlog{dyn_deaths}
\cgeneed{22\baselineskip}
\stlog{dyn_health}
Health services per head rise, and the death rate falls by 0.022 points per
1,000 in period 1. Period 0 has the base death rate, so the population grows
at \cmd{population()}, 2 per cent, in period 1. From period 2, the lower death
rate of the period before adds to population growth: the population index is
4.042 per cent above the base in period 2, against 4.040 at a constant 2 per
cent ($1.02^2 = 1.0404$). Labour supply grows by $0.3/0.55$ times the fall in
the death rate in addition to its rate of 2 per cent, and is 4.041 per cent
above the base in period 2. Because rising health services per head lower the
death rate in a baseline as well, scenarios should be compared with a
baseline run with the same settings.

\cgesubsec{Example 4: investment with perfect foresight}

In a recursive run investment cannot respond to a policy before it takes
effect. With perfect foresight it can. Savings adjust to investment
(\cmd{savinv(investment)}), and investment follows Tobin's $q$, the value of a
unit of capital that earns the real rental rate $\V{RK} = \V{W}_{cap}/\V{CPI}$
and depreciates at 5 per cent, at an interest rate of 5 per cent:
$Q_t = (\V{RK}_{t+1} + 0.95\,Q_{t+1})/1.05$. Its base value, $1/0.10 = 10$,
is the steady state of the base. The swap fixes $\V{IADJ}/(Q/Q^0)^2$, so that
the investment factor $\V{IADJ}$ moves with the square of $q$:
\stlog{dyn_fs_setup}
A baseline without shocks, and a run in which every import tariff is removed
from period 4, known from period 0; capital accumulates from investment, at a
gross rental rate of 10 per cent, and labour grows at 2 per cent:
\stlog{dyn_fs_run}
The stacked system has the 96 equations of each of the 11 periods and the
capital stock of each period after the first. Newton's method starts from a
recursive run in which each period expects the later ones to repeat the period
before, and converges in a few iterations. The investment factor against the
baseline:
\stlog{dyn_fs_compare}
In both runs the investment factor is above 1 from period 0: labour grows at 2
per cent a period, which raises the future return to capital, and $q$ prices
it at once. The tariff cut adds 2.9 per cent to investment in period 0, four
periods before the tariffs go, and 4.0 per cent by period 3: the higher return
to capital from period 4 is known in advance, and capital is built for it.
From the last period the leads repeat its values (\cmd{terminal(constant)}),
so investment is the same in periods 9 and 10. A recursive run cannot show the
early response: in it the cut changes investment only from period 4.

\cgesubsec{Capital}

Capital factors are listed in \cmd{capital()}. New capital comes from the
previous period's real investment at base prices, shared among capital
factors in proportion to their base income. When capital is mobile, it adds to
the factor's supply. When it is specific to each activity
(\manref{cge_closure}, \cmd{factors(cap:specific)}), new capital is allocated
among the activities that use the factor in proportion to their capital,
tilted towards those whose rental rate is above the average by
\cmd{sensitivity()}; an activity's share cannot fall below zero. The rental
rate of \cmd{rental()} converts investment, measured at base prices, into the
factor's units, whose base quantity is the factor's base income: with a
single capital factor and a rate of 0.15, one unit of real investment adds
0.15 units of the factor, and the base stock is base capital income divided
by 0.15.

\cgesubsec{Labour force and unemployment}

A factor with a labour force (\manref{cge_calibrate}, \cmd{unemployment()} or
\cmd{minwage()}) has its labour force grow at the factor's growth rate. Under
unemployment at a fixed wage, a wage curve or a minimum wage, employment is
then determined by labour demand, and the unemployment rate changes
(\chapref{chap:labour}).

\cgesubsec{Public debt}

When the model has public debt (\manref{cge_calibrate}, \cmd{debt()}), the
deficit of each period is borrowed in the next, the share given by
\cmd{foreignborrowing()} of \manref{cge_calibrate} abroad, and interest is paid
on the new stocks; a surplus repays debt in the same shares
(\secref{sec:debt}). A program in \cmd{shocks()} may change the stocks and the
interest rates, for example \cmd{cge_shock fdebt, scale(0.5)} in one period for
debt relief.

\cgesubsec{Shocks and the model in memory}

A run starts from the base: shocks set with \manref{cge_shock} before the run
are not used, and a note is displayed when fixed values have been changed.
Constraints set by \manref{cge_constraint} are kept. Other fixed values,
including the numeraire, keep their base values unless a program in
\cmd{shocks()} changes them. When the run ends, or stops with an error, the
parameters and fixed values that the model had before it are restored. A
period that does not converge stops the run with error 430; the periods
solved before it remain stored.

\cgesubsec{Lags in user equations}

From version 0.9.2, user equations (\manref{cge_equation}) may refer to the
values of variables in the periods before, \cmd{lag(GL)} or
\cmd{lag(QA(act_man), 2)}, for up to four periods; before the first period
they are the base values. Stocks and partial adjustment that the run does not
update itself, such as a fund that accumulates or spending that follows
revenue with a delay, can so be written in the model's own terms. Each period
stores the lagged values it used.

\cgesubsec{Results}

\manref{cge_report} with \cmd{run()} reports every period of a run, and
\manref{cge_dynresults} returns every variable of every period as a dataset,
with the differences from a baseline run. A period can be passed to
\manref{cge_microsim} like any other simulation; its changes are measured from
the base. \mbox{\manref{cge_gams}} with \cmd{dynamic()} writes a run as a GAMS
program that solves the periods in order.

\cgesec{Stored results}

\cmd{cge_dynamic} stores the following in \cmd{r()}:
\begin{storedres}
\resgroup{Macros}
\res{r(name)} & name of the run\\
\res{r(sims)} & simulations of the run, in period order\\
\res{r(foresight)} & \cmd{foresight} for a run with perfect foresight\\
\end{storedres}

\cgesec{Methods and formulas}

Let periods $t = 0, 1, \dots, T$ be the steps of \cmd{periods()}, with $t=0$
the base. Before period $t \ge 1$, from the solution of period $t-1$:

\textit{Capital.} For each capital factor $f \in K$, with depreciation rate
$\delta_f$ and base gross rental rate $\rho_f$, new capital is
\begin{equation*}
N_{f,t} = \rho_f\,\omega_f\,\V{INV}_{t-1}, \qquad
\V{INV}_{t-1} = \sum_c \V{PQ}^0_c\,\V{QINV}_{c,t-1}, \qquad
\omega_f = \frac{\V{YF}^0_f}{\sum_{f' \in K} \V{YF}^0_{f'}} .
\end{equation*}
For mobile capital, $\V{QFS}_{f,t} = (1-\delta_f)\,\V{QFS}_{f,t-1} + N_{f,t}$.
For capital specific to each activity,
\begin{equation*}
\V{QF}_{fa,t} = (1-\delta_f)\,\V{QF}_{fa,t-1} + N_{f,t}\,\eta_{fa,t}, \qquad
\eta_{fa,t} = \frac{\max\bigl[0,\; s_{fa}\bigl(1 + \beta(R_{fa}/\bar R_f - 1)\bigr)\bigr]}
{\sum_{a'} \max\bigl[0,\; s_{fa'}\bigl(1 + \beta(R_{fa'}/\bar R_f - 1)\bigr)\bigr]},
\end{equation*}
over the activities that use the factor, with
\begin{equation*}
s_{fa} = \frac{\V{QF}_{fa,t-1}}{\sum_{a'} \V{QF}_{fa',t-1}}, \qquad
R_{fa} = \V{W}_{f,t-1}\,\V{WFDIST}_{fa,t-1}, \qquad
\bar R_f = \sum_a s_{fa}\,R_{fa},
\end{equation*}
and $\beta$ the value of \cmd{sensitivity()}.

\textit{Population and other factors.} With a death rate, base rate $m^0$ and
elasticity $\epsilon^m$, the death rate of period $t-1$, with the parameters
and health services of that period, is
\begin{equation*}
\V{MR}_{t-1} = \mu^m \left(\frac{\V{HS}_{t-1}}{\V{HS}^0_{t-1}}\right)^{\epsilon^m},
\end{equation*}
where $\mu^m = m^0$ unless it is shocked, and
\begin{equation*}
g_{P,t} = g_P + \bigl(m^0 - \V{MR}_{t-1}\bigr), \qquad
g_{f,t} = g_f + d_f\,\frac{\theta^D}{\alpha^W}\bigl(m^0 - \V{MR}_{t-1}\bigr),
\end{equation*}
where $g_P$ is \cmd{population()}, $g_f$ the rate of factor $f$ in
\cmd{growth()}, $d_f = 1$ for the factors in \cmd{deathfactors()} and 0
otherwise, $\theta^D$ is \cmd{deathshare()} and $\alpha^W$ \cmd{workingage()}.
Without a death rate, $g_{P,t} = g_P$ and $g_{f,t} = g_f$. A factor that is
not capital grows as
\begin{equation*}
\V{QFS}_{f,t} = (1 + g_{f,t})\,\V{QFS}_{f,t-1} \quad\text{(mobile)}, \qquad
\V{QF}_{fa,t} = (1 + g_{f,t})\,\V{QF}_{fa,t-1} \quad\text{(specific)},
\end{equation*}
and its labour force, if it has one, as $\V{qlf}_{f,t} = (1+g_{f,t})\,\V{qlf}_{f,t-1}$.
The subsistence quantities $\gamma_{ch}$ and $\gamma^h_{ah}$, the base of the
health link $\V{HS}^0$ and the population index grow at $g_{P,t}$.

\textit{Government and foreign flows.} With $g_G$ the rate of
\cmd{govgrowth()} and $g_W$ that of \cmd{foreigngrowth()},
\begin{align*}
&\overline{\V{qg}}_c,\ \V{trnsfr}_{d,gov},\ \V{trnsfr}_{gov,gov},\ \V{trnsfr}_{shi,gov}
  && \text{grow at } g_G,\\
&\V{trnsfr}_{d,row},\ \V{trnsfr}_{f,row},\ \V{trnsfr}_{gov,row},\ \V{trnsfr}_{row,gov},\
  \V{trnsfr}_{shi,row} && \text{grow at } g_W,
\end{align*}
and so does $\V{FSAV}$ when the closure fixes it.

\textit{Public debt.} With debt and a share $\phi$ borrowed abroad,
\begin{equation*}
\V{DD}_t = \V{DD}_{t-1} + (1-\phi)\,(-\V{GSAV}_{t-1}), \qquad
\V{DF}_t = \V{DF}_{t-1} + \phi\,(-\V{GSAV}_{t-1})/\V{EXR}_{t-1} .
\end{equation*}

\textit{Shocks and efficiency.} The program in \cmd{shocks()} is called, and
efficiency is set to
\begin{equation*}
\V{ad}_{a,t} = \lambda_{a,t}\,\V{ad}_{a,t-1},
\end{equation*}
with $\lambda_{a,t} = 1 + g^{ad}_a$ under \cmd{tfp()}, the uniform factor
$\lambda_t$ of the baseline's period $t$ under \cmd{baseline()}, and $1$
otherwise. Under \cmd{gdpgrowth()}, $\lambda_t$ is uniform and solves
\begin{equation*}
\V{VA}_t(\lambda_t) = (1 + g^Y_t)\,\V{VA}_{t-1}, \qquad
\V{VA}_t = \sum_a \V{PVA}^0_a\,\V{QVA}_{a,t},
\end{equation*}
real value added at base prices, less commodities bought within nests of
value added (\secref{sec:nests}). It is found by secant iterations that
start from $\lambda = 1$ and $\lambda = 1 + g^Y_t$, each of which solves the
model, until $\lvert \V{VA}_t / ((1+g^Y_t)\,\V{VA}_{t-1}) - 1 \rvert < 10^{-13}$,
in at most 50 iterations; otherwise the run stops with error 430.

Period $t$ is then solved from the solution of period $t-1$ by the solver of
\manref{cge_solve}, without continuation. \Chapref{chap:dynamics} states the
dynamics in full, and \secref{sec:population} the population and labour
updates.

\cgesec{References}

Thurlow, J. 2004. \textit{A dynamic computable general equilibrium (CGE) model
for South Africa: Extending the static IFPRI model}. Working Paper 1-2004,
Trade and Industrial Policy Strategies.

\cgesec{Also see}

\begin{alsosee}
\seeentry{cge_dynresults}{Results of a dynamic run as a dataset}
\seeentry{cge_report}{Standard report tables for one or more simulations}
\seeentry{cge_shock}{Change a parameter or a fixed variable}
\seeentry{cge_closure}{Choose the model closure}
\seeentry{cge_calibrate}{Calibrate the model to the SAM}
\seeentry{cge_gams}{Write the model as a GAMS program, and solve it in GAMS}
\seechapter{chap:dynamics}{Recursive dynamics}
\end{alsosee}
